% 2.1.9
% Ficheiro incluído por main.tex

\subsection[Acesso à SRII como utilizador do domínio]{Acesso à SRII como utilizador do domínio}

\subsubsection[Aluno]{Aluno}

\passo{Passo 1 - Obter o ticket na sessão do aluno}

Continuando na sessão do aluno (2.1.8):

\begin{lstlisting}[style=bashstyle]
kinit aluno.teste@ISEC.LOCAL
klist
\end{lstlisting}

\emph{Verificação esperada:} \texttt{klist} mostra \texttt{krbtgt/\allowbreak{}ISEC.LOCAL@ISEC.LOCAL}.

\passo{Passo 2 - Listar a SRII com o ticket}

\begin{lstlisting}[style=bashstyle]
smbclient //SR2-SRV.isec.local/SRII --use-kerberos=required -c 'ls'
\end{lstlisting}

\emph{Verificação esperada:} a listagem da SRII aparece.

\passo{Passo 3 - Testar a escrita}

\begin{lstlisting}[style=bashstyle]
smbclient //SR2-SRV.isec.local/SRII --use-kerberos=required -c 'put /tmp/teste-escrita.txt aluno-escrita2.txt'
klist
exit
\end{lstlisting}

\emph{Verificação esperada:} o \texttt{put} falha com \texttt{NT\_STATUS\_ACCESS\_DENIED}; o \texttt{klist} mostra dois tickets: \texttt{krbtgt/\allowbreak{}ISEC.LOCAL} e \texttt{cifs/\allowbreak{}SR2-SRV.isec.local}.

\subsubsection[Professor]{Professor}

\passo{Passo 1 - Entrar como professor}

\begin{lstlisting}[style=bashstyle]
su - professor.teste@isec.local
\end{lstlisting}

\passo{Passo 2 - Obter o ticket na sessão}

\begin{lstlisting}[style=bashstyle]
kinit professor.teste@ISEC.LOCAL
\end{lstlisting}

\passo{Passo 3 - Ler e escrever na SRII}

\begin{lstlisting}[style=bashstyle]
smbclient //SR2-SRV.isec.local/SRII --use-kerberos=required -c 'ls'
smbclient //SR2-SRV.isec.local/SRII --use-kerberos=required -c 'put /tmp/teste-escrita.txt professor-escrita2.txt'
smbclient //SR2-SRV.isec.local/SRII --use-kerberos=required -c 'ls'
klist
exit
\end{lstlisting}

\emph{Verificação esperada:} o professor lista e escreve; \texttt{professor-escrita2.txt} aparece na segunda listagem.

